\section{Strongly Connected Components (Kosaraju's Algorithm)}
\label{sec:graph-kosaraju}

\problemstatement{Given a directed graph, find its \textbf{Strongly
Connected Components} (SCCs) -- maximal groups of nodes where every
node can reach every other node in the same group via directed
edges.}

\begin{patternbox}
\emph{``Group nodes that can all reach each other''} is solved by
Kosaraju's algorithm: run DFS once to compute a finish-time order
(exactly the Graph Topological Sort chapter's finish-time idea), then
run DFS \textbf{again} on the graph with every edge \textbf{reversed},
processing nodes in \textbf{decreasing} finish-time order -- each DFS
launched in this second pass visits exactly one complete SCC.
\end{patternbox}

\subsection*{Why reversing edges and using finish-time order works}

Reversing every edge preserves SCCs exactly (if $u$ and $v$ can reach
each other in the original graph, they can still reach each other with
every edge flipped) but \textbf{destroys} inter-SCC reachability in a
specific, useful direction: collapsing each SCC into a single
``super-node'' produces a DAG (the \textbf{condensation} of the
graph), and processing nodes in decreasing finish-time order from the
\emph{first} DFS guarantees the first unvisited node encountered in
the second pass always belongs to a \textbf{source} SCC (no incoming
edges from any other not-yet-visited SCC) of the \emph{original}
condensation -- which is a \textbf{sink} SCC once edges are reversed.
Starting DFS there in the reversed graph cannot leak into any other
SCC, since reversed edges now only point \emph{into} this SCC from
elsewhere, not out of it.

\subsection*{Algorithm}

\begin{enumerate}
  \item Run DFS over the original graph from every unvisited node,
        pushing each node onto a stack when it finishes (identical to
        topological sort's DFS approach).
  \item Build the transposed graph (every edge reversed).
  \item Pop the stack one node at a time; for every still-unvisited
        node, run DFS on the \textbf{transposed} graph -- everything
        reached in this one DFS call is exactly one SCC.
\end{enumerate}

\subsection*{C++ implementation}

\begin{lstlisting}
#include <bits/stdc++.h>
using namespace std;

void fillOrder(int node, vector<vector<int>>& adj, vector<bool>& visited, stack<int>& order) {
    visited[node] = true;
    for (int neighbor : adj[node]) {
        if (!visited[neighbor]) fillOrder(neighbor, adj, visited, order);
    }
    order.push(node);
}

void collectSCC(int node, vector<vector<int>>& transposed, vector<bool>& visited,
                 vector<int>& component) {
    visited[node] = true;
    component.push_back(node);
    for (int neighbor : transposed[node]) {
        if (!visited[neighbor]) collectSCC(neighbor, transposed, visited, component);
    }
}

vector<vector<int>> kosaraju(int n, vector<vector<int>>& adj) {
    vector<bool> visited(n, false);
    stack<int> order;
    for (int i = 0; i < n; i++) {
        if (!visited[i]) fillOrder(i, adj, visited, order);
    }

    vector<vector<int>> transposed(n);
    for (int u = 0; u < n; u++) {
        for (int v : adj[u]) transposed[v].push_back(u);
    }

    fill(visited.begin(), visited.end(), false);
    vector<vector<int>> sccs;
    while (!order.empty()) {
        int node = order.top();
        order.pop();
        if (!visited[node]) {
            vector<int> component;
            collectSCC(node, transposed, visited, component);
            sccs.push_back(component);
        }
    }
    return sccs;
}

int main() {
    int n = 5;
    vector<vector<int>> adj(n);
    adj[0] = {2}; adj[2] = {1}; adj[1] = {0}; adj[0].push_back(3); adj[3] = {4};

    for (auto& scc : kosaraju(n, adj)) {
        for (int x : scc) cout << x << " ";
        cout << endl;
    }
    // 0 1 2   (one SCC)
    // 3        (its own SCC)
    // 4        (its own SCC)
    return 0;
}
\end{lstlisting}

\subsection*{Dry run}

\dryrun{Take edges $0{\to}2,2{\to}1,1{\to}0,0{\to}3,3{\to}4$.}

First DFS (finish-time order, starting from $0$): visits $0\to2\to1$
(back to $0$, already visiting, skip), $1$ finishes first (push),
$2$ finishes (push), then $0$ continues to $3\to4$, $4$ finishes
(push), $3$ finishes (push), $0$ finishes (push). Stack
(bottom-to-top): $1,2,4,3,0$. Transposed graph: $2{\to}0,1{\to}2,0{\to}1,
3{\to}0,4{\to}3$. Second DFS, popping stack top-to-bottom: pop $0$
first -- DFS on transposed graph from $0$: $0\to1\to2\to0$ (visited,
stop) -- component $\{0,1,2\}$. Pop $3$ (next unvisited): DFS from $3$
on transposed graph: $3\to0$ (visited, stop) -- component $\{3\}$.
Pop $4$: component $\{4\}$. Three SCCs: $\{0,1,2\}$, $\{3\}$, $\{4\}$
-- correctly identifying the $3$-cycle as one SCC and the two
remaining nodes (no return path) as singleton SCCs.

\begin{complexitybox}
\textbf{Time Complexity.} $O(V+E)$ -- two full DFS passes plus
building the transpose, each linear in the graph size.
\textbf{Space Complexity.} $O(V+E)$ for the transposed adjacency list,
plus $O(V)$ for the stack and visited arrays.
\end{complexitybox}

\begin{takeawaybox}
Kosaraju's algorithm is this book's clearest demonstration that
\textbf{reversing} a directed graph's edges is sometimes itself the
key algorithmic move -- not just bookkeeping, but the step that turns
an otherwise hard global structure question into two simple, linear
DFS passes.
\end{takeawaybox}
